\begin{tcolorbox}[colback=red!3!white,colframe=red!50!black,boxrule=0.8pt,title={Korrekturhinweise},fonttitle=\bfseries]
\textbf{Hinweis zur Rechenprüfung (29.~Sept.~2026).} Einzelne Rechnungen:
\begin{itemize}
\item $f = 1/(4\xi) = 1875$, nicht $7500$; richtig ist $f = 1/\xi = 30000/4 = 7500$.
\item $7500 = 2^2\cdot3\cdot5^4$ hat $30$ positive Teiler, nicht $36$.
\item $G = \frac{\xi}{2}\cdot10^{-6} = 6{,}667\times10^{-11}$, also $-0{,}11\,\%$, nicht $0{,}01\,\%$.
\item Die Quarkformeln ergeben $u$: $60{,}5$, $d$: $80{,}6$, $s$: $156$~MeV, $c$: $2{,}73$, $b$: $4{,}42$~GeV statt der angegebenen Werte; kein gemeinsamer Anker erklärt die Abweichungen. Die $t$-Formel stimmt auf etwa $1\,\%$.
\end{itemize} Der übrige Text ist unverändert.
\end{tcolorbox}

\section*{Abstract}
		In der vorliegenden Arbeit wird die fundamentale Architektur der Raumzeit im Rahmen der \textbf{Fundamental Fractal Geometric Field Theory (FFGFT)} – intern als T0-Modell bezeichnet – neu interpretiert. Das zentrale Paradigma besteht im Übergang von einer punktförmigen zu einer rein geometrischen Beschreibung des Vakuums als vierdimensionaler \textbf{Hirnwindungs-Torus}.
		
		\textbf{Geometrischer Aufbau:} Die Theorie gründet auf der fraktal-geometrischen Grundstruktur mit dem Parameter $\xi \approx (4/3)\times 10^{-4}$ und der dichtesten lokalen Kugelpackung durch reguläre \textbf{Tetraeder}. Diese tetraedrische Basis bildet das stabile Fundament für die niedrigen Generationen (Elektron, Myon, Proton/Neutron) sowie die lokale 3D-Kristallstruktur des Torsos. Darauf aufbauend entsteht durch fraktale Verzweigung und pentagonale Symmetriebrechung der ideale sub-Planck-Faktor
		\begin{equation*}
			f = 7500,
		\end{equation*}
		der eine exakt 7500-fache Verkleinerung gegenüber der konventionellen Planck-Skala ($t_0$) darstellt und direkt aus der geometrischen Windungsdichte $30000/4$ folgt.
		
		\textbf{g-2-Anomalie:} Ein Kernstück der Arbeit ist die transparente geometrische Herleitung der anomalen magnetischen Momente der Leptonen. Während das Standardmodell auf zahlreiche störungstheoretische Terme angewiesen ist, ergibt sich in der FFGFT die Elektron-Anomalie direkt aus der Basiswindung (tetraedrische Projektion). Die Myon- und Tau-Anomalien entstehen durch fraktale Verzweigungen mit den Hausdorff-Dimensionen $p \approx 5/3$ bzw. $4/3$. Mit dem idealen Wert $f = 7500$ erreichen die rein geometrischen Vorhersagen eine Genauigkeit von etwa 2\,\%. Durch Rekonstruktion des Projektionsfaktors $k_\text{geom}$ sinkt die Abweichung beim Myon auf unter 0{,}2\,\%. Die präziseste, $k_\text{geom}$-unabhängige Vorhersage für die Tau-Anomalie lautet
		\begin{equation*}
			a_\tau \approx 1{,}282 \times 10^{-3},
		\end{equation*}
		die ausschließlich aus dem exakten Verhältnis $f^{1/3} - 1$ folgt.
		
		\textbf{Geometrische Verhältnismäßigkeit:} Alle physikalischen Basisgrößen (Konstanten, Massen, Kopplungen) werden auf feste geometrische Verhältnisse zurückgeführt; die Theorie kommt mit dem einen Parameter $\xi$ aus (plus motivierter ganzzahliger Strukturwahlen), während SI-Anker der Umrechnung dienen. Die T0-Theorie bietet so eine transparente geometrische Beschreibung und liefert konkrete, experimentell überprüfbare Vorhersagen – insbesondere für die Tau-Anomalie als klaren Test bei Belle II.

	\medskip\noindent\textbf{Statusmarker.} Aussagen mit \textbf{[S]} sind \emph{spekulativ}: ein Hinweis oder eine Vermutung, die im Korpus noch nicht hergeleitet oder numerisch geprüft ist.

	
	
	
	\section{Einleitung: Das geometrische Paradigma}
	
	\subsection{Offene Fragen der modernen Physik}
	
	Das 21. Jahrhundert steht vor einer grundlegenden Frage: Während das Standardmodell der Teilchenphysik experimentelle Daten mit sehr hoher Präzision beschreibt, enthält es doch ~19 freie Parameter, die nicht aus Prinzipien abgeleitet werden können, sondern empirisch angepasst werden müssen. Zudem sagt dieses Modell keine Werte für fundamentale Konstanten wie die Feinstrukturkonstante \(\alpha\), die Massen von Elektron oder Proton, oder die Stärke der Gravitation voraus.
	
	Gleichzeitig gibt es Hinweise auf Phänomene, die über das Standardmodell hinausweisen könnten: Die beobachtete Beschleunigung der kosmischen Expansion (Dunkle Energie), die Anomalien in den Rotationskurven von Galaxien (Dunkle Materie), und die präzisen Messungen der anomalen magnetischen Momente von Leptonen zeigen alle Diskrepanzen zur etablierten Theorie.
	
	Die T0-Theorie verfolgt einen anderen Ansatz: Statt neue Teilchen oder Felder zu postulieren, geht sie von einer fundamentalen geometrischen Struktur der Raumzeit selbst aus.
	
	\subsection{Die Grundidee: Raumzeit als Torsionskristall}
	
	Die zentrale These der T0-Theorie lässt sich in einem Satz zusammenfassen:
	
	\begin{center}
		\textbf{Das Universum ist ein statischer 4-dimensionaler Torsionskristall, dessen diskrete Sub-Planck-Struktur alle beobachtbaren physikalischen Phänomene erzeugt.}
	\end{center}
	
	Was bedeutet das konkret?
	
	\begin{enumerate}
		\item \textbf{Statisch \textbf{[S]}:} Das Universum expandiert nicht im herkömmlichen Sinne. Die beobachtete Rotverschiebung entsteht durch geometrische Wegverlängerung im Torsionsgitter. (Kosmologische Aussage; der heutige Korpus klammert den kosmischen Sektor bewusst aus, weil auch das Standardmodell ihn nicht herleitet, P39.)
		\item \textbf{4-dimensional:} Neben den drei räumlichen Dimensionen existiert eine vierte, die nicht mit der Zeit identisch ist, sondern eine zusätzliche räumliche Dimension darstellt, die in unserem Erfahrungsraum \enquote{aufgerollt} ist.
		\item \textbf{Torsionskristall:} Raumzeit ist nicht kontinuierlich, sondern besitzt auf der Sub-Planck-Skala eine diskrete, kristalline Struktur. Die \enquote{Torsion} beschreibt die Windungen und Verdrillungen dieser Kristallstruktur.
		\item \textbf{Sub-Planck-Struktur:} Die fundamentale Längenskala ist nicht die Planck-Länge \(\ell_P = 1{,}616 \times 10^{-35}\,\text{m}\), sondern eine um den Faktor \(f = 7491{,}91\) kleinere Skala.
	\end{enumerate}
	
	In diesem Bild sind \textbf{Teilchen keine punktförmigen Objekte}, sondern stehende Wellen (Resonanzen) im Torsionskristall. \textbf{Kräfte} sind nicht Austausch virtueller Teilchen, sondern geometrische Kopplungen zwischen verschiedenen Torsionsmoden. \textbf{Massen} sind keine intrinsischen Eigenschaften, sondern Frequenzen dieser Resonanzen.
	
	\section{Die fundamentale Herleitung: Von der Geometrie zum Zahlenwert}
	\subsection{Der narrative Ausgangspunkt: Warum 30000?}
	Die Herleitung beginnt mit einer scheinbar willkürlichen Zahl: \(30000\). Sie ist eine motivierte Setzung: Sie soll die Struktur der 4-dimensionalen Raumzeit kodieren.
	Stellen Sie sich vor: Wir leben in einer Welt mit \textbf{drei} erfahrbaren Raumdimensionen. Doch auf fundamentalster Ebene existiert eine \textbf{vierte} Dimension, die nicht direkt zugänglich ist, sondern nur indirekt durch ihre geometrischen Effekte spürbar wird. Diese vierte Dimension ist \enquote{kompaktifiziert} -- sie ist auf kleinsten Skalen aufgerollt.
	Die Zahl \(30000\) entsteht aus der Wechselwirkung zwischen diesen vier Dimensionen:
	\begin{itemize}
		\item Die \textbf{3} steht für die drei erfahrbaren Raumdimensionen.
		\item Die \textbf{4} steht für die volle, vierdimensionale Realität.
		\item Die \textbf{000} (also Faktor 1000) beschreibt die Skalenhierarchie zwischen der fundamentalen und der beobachtbaren Ebene.
	\end{itemize}
	Konkret definieren wir:
	\begin{equation}
		\boxed{\xi = \frac{4}{30000} = 1{,}333\overline{3} \times 10^{-4}}
	\end{equation}
	Diese Zahl \(\xi\) ist der \textbf{fundamentale Korrekturparameter}. Sie beschreibt, wie stark die reale 4D-Raumzeit von einer idealen 3D-Geometrie abweicht. Physikalisch interpretiert: \(\xi\) ist die \enquote{Torsionsspannung} -- die winzige Verwindung, die das Raumzeit-Gitter von einer perfekten Struktur unterscheidet.
	\subsection{Die ideale Ankerzahl: Warum 7500?}
	Aus \(\xi\) ergibt sich die ideale Ankerzahl:
	\begin{equation}
		\boxed{f = \frac{1}{4\xi} = \frac{30000}{4} = 7500}
	\end{equation}
	Dies ist die Zahl, die als idealer Sub-Planck-Faktor bezeichnet wird: Die \textbf{ideale Ankerzahl} des Kristallgitters.
	\textbf{Warum ist 7500 so speziell?} Schauen wir uns die Primfaktorzerlegung an:
	\begin{equation}
		7500 = 2^2 \times 3 \times 5^4 = 4 \times 3 \times 625
	\end{equation}
	Die Zahl hat eine reiche Teilerstruktur:
	\begin{itemize}
		\item Sie hat \textbf{36 positive Teiler} -- günstig für eine symmetrische Gitterstruktur.
		\item Sie kombiniert die ersten drei Primzahlen (2, 3, 5).
		\item Der Faktor \(5^4 = 625\) verweist auf die pentagonale Symmetrie des Kristalls (5) in vier Dimensionen (Exponent 4).
		\item Die Zahl ist durch zahlreiche Faktoren teilbar -- eine günstige Basis für Resonanzen.
	\end{itemize}
	In der Kristallographie bezeichnet man Strukturen mit vielen Teilern als \enquote{hochsymmetrisch} -- das passt zu einer fundamentalen Raumzeitstruktur.
	\subsection{Die Symmetriebrechung: Die Rolle des goldenen Schnitts}
	Ein perfekter, idealer Kristall wäre vollkommen symmetrisch. Doch unsere Welt zeigt Symmetriebrechungen auf allen Ebenen:
	\begin{itemize}
		\item Materie dominiert über Antimaterie
		\item Die schwache Wechselwirkung verletzt die Paritätssymmetrie
		\item Das Neutron ist schwerer als das Proton
		\item Die drei Generationen der Leptonen haben unterschiedliche Massen
	\end{itemize}
	Die T0-Theorie schlägt \textbf{[S]} für diese Symmetriebrechungen einen gemeinsamen geometrischen Ursprung vor: die pentagonale Symmetrie des Kristalls, verkörpert durch den \textbf{goldenen Schnitt} \(\varphi\).
	Der goldene Schnitt \(\varphi = (1+\sqrt{5})/2 = 1{,}618033989\ldots\) ist die irrationale Zahl, die die pentagonale Symmetrie beschreibt. In einem perfekten Fünfeck taucht \(\varphi\) überall auf: Das Verhältnis von Diagonale zu Seite ist genau \(\varphi\).
	Warum ausgerechnet pentagonale Symmetrie? Aus mathematischen Gründen (kristallographische Beschränkung) ist die pentagonale Symmetrie die erste, die in der Ebene \textbf{nicht periodisch parkettieren} kann. Dies führt zu \enquote{Quasikristallen} -- Strukturen, die geordnet, aber nicht periodisch sind. Genau eine solche quasikristalline Struktur postuliert die T0-Theorie für die Sub-Planck-Skala.
	Die Symmetriebrechung wird in der Theorie nicht durch eine direkte Subtraktion von \(5\varphi\) von der idealen Ankerzahl 7500 quantifiziert. Stattdessen ist sie in den \textbf{ca. 2\,\% Abweichungen} verborgen, die in den Berechnungen der anomalen magnetischen Momente (g-2-Anomalien) auftreten. Diese Abweichung entsteht durch die pentagonale Projektion in den geometrischen Faktor \(k_\text{geom}\):
	\begin{equation}
		k_\text{geom} = \frac{2}{\sqrt{\varphi}} \times \sqrt{2} \approx 2{,}22357,
	\end{equation}
	der die 4D-Torsion auf die 3D-Welt projiziert. Die rekonstruierte Version aus experimentellen Daten weicht um etwa 2\,\% ab (\(k_\text{geom}^\text{rek} \approx 2{,}26955\)), was die eigentliche Symmetriebrechung widerspiegelt -- eine leichte Verzerrung durch die pentagonale Geometrie, die die perfekte Symmetrie bricht, ohne den idealen Wert \(f = 7500\) zu verändern.
	\subsection{Der reale Sub-Planck-Faktor: \(f = 7500\)}
	Nun setzen wir alles zusammen: Der ideale Kristall bleibt erhalten, die Symmetriebrechung wirkt sich nur in den Projektionsfaktoren aus:
	\begin{equation}
		\boxed{f = 7500}
	\end{equation}
	Dies ist eine \textbf{zentrale Zahl der T0-Theorie}. Sie erscheint in fast allen Formeln und beschreibt:
	\begin{itemize}
		\item Die Anzahl der Sub-Planck-Zellen pro Planck-Länge
		\item Die Dichte des Torsionsgitters
		\item Die Grundfrequenz aller geometrischen Resonanzen
	\end{itemize}
	\section{Stufe 1: Von der Geometrie zur Energie -- das Higgs-Feld}
	
	\subsection{Die Planck-Skala als natürliche Referenz}
	
	In der theoretischen Physik gibt es eine natürliche Skala für Masse, Länge und Zeit: die Planck-Skala. Diese ergibt sich aus einer Kombination der fundamentalen Konstanten:
	\begin{align}
		m_P &= \sqrt{\frac{\hbar c}{G}} = 1{,}220910 \times 10^{19}\,\text{GeV}/c^2 \\
		\ell_P &= \sqrt{\frac{\hbar G}{c^3}} = 1{,}616255 \times 10^{-35}\,\text{m} \\
		t_P &= \sqrt{\frac{\hbar G}{c^5}} = 5{,}391247 \times 10^{-44}\,\text{s}
	\end{align}
	
	Diese Größen markieren die Skala, bei der Quanteneffekte der Gravitation wichtig werden. In der herkömmlichen Physik bleibt unklar, warum die beobachtbaren Teilchenmassen so viel kleiner sind als die Planck-Masse (das Hierarchieproblem).
	
	In der T0-Theorie erhält die Planck-Skala eine geometrische Interpretation: Sie ist die \textbf{Gitterschwingungsfrequenz} des fundamentalen Kristalls. Die Planck-Masse ist die Energie, die benötigt wird, um eine einzelne Gitterzelle maximal anzuregen.
	
	\subsection{Die 4D-Energiedichte: Verdünnung über vier Dimensionen}
	
	Die Grundannahme dieses Abschnitts ist: Die Planck-Energie wird nicht auf einer einzigen Zelle konzentriert, sondern verteilt sich über das vierdimensionale Gitter. Warum vier Dimensionen? Weil jede der vier Raumdimensionen des Torsionskristalls zur Energiedichte beiträgt.
	
	Mathematisch bedeutet dies:
	\begin{equation}
		\boxed{\rho_{4D} = \frac{m_{\text{Planck}}}{f^4}}
	\end{equation}
	
	\textbf{Narrative Erklärung:} Stellen Sie sich einen perfekten Würfel vor, dessen Kantenlänge \(f\) Zellen beträgt. In drei Dimensionen enthält dieser Würfel \(f^3\) Zellen. In vier Dimensionen enthält der Hyperwürfel \(f^4\) Zellen. Die Planck-Energie, die ursprünglich auf einer einzelnen Zelle konzentriert war, verteilt sich nun gleichmäßig über alle \(f^4\) Zellen des vierdimensionalen Hyperwürfels.
	
	Rechnen wir nach:
	\begin{equation}
		f^4 = 7491{,}91^4 \approx 3{,}155 \times 10^{15}
	\end{equation}
	
	Die 4D-Energiedichte ist also um den Faktor \(3{,}155 \times 10^{15}\) kleiner als die Planck-Masse:
	\begin{equation}
		\rho_{4D} = \frac{1{,}220910 \times 10^{19}\,\text{GeV}}{3{,}155 \times 10^{15}} \approx 3{,}869 \times 10^{3}\,\text{GeV}
	\end{equation}
	
	Wir erhalten eine Energiedichte von etwa 3869 GeV. Dies ist immer noch viel höher als die beobachtbaren Energieskalen, aber wir sind auf dem richtigen Weg.
	
	\subsection{Projektion auf 3D: Der Halbraum-Effekt}
	
	Wir leben in einer dreidimensionalen Welt. Die vierte Dimension ist für uns nicht direkt zugänglich. Wie kommt die Energiedichte aus der vierten Dimension in unsere dreidimensionale Erfahrungswelt?
	
	Dies geschieht durch \textbf{geometrische Projektion}. Stellen Sie sich eine 4D-Kugel (eine 3-Sphäre) vor, die in unsere 3D-Welt projiziert wird. Die Projektion einer vollen 4D-Kugel auf den 3D-Halbraum erfolgt durch Division durch \(\pi/2\).
	
	Warum gerade \(\pi/2\)? Betrachten wir den einfacheren 2D-Fall: Die Projektion eines Halbkreises (Winkel \(\pi\)) auf eine Gerade ergibt einen Faktor \(\pi/2\). Analog ist die Projektion einer 3-Sphäre (Oberfläche: \(2\pi^2\)) auf den 3D-Halbraum durch \(\pi/2\) gegeben.
	
	\subsection{Skalierung auf die elektroschwache Skala: Der Faktor 1/10}
	
	Die nach Projektion erhaltene 3D-Energiedichte muss noch auf die elektroschwache Skala skaliert werden. Der Übergang von der fundamentalen geometrischen Skala zur elektroschwachen Skala erfordert eine weitere Skalierung um Faktor 1/10.
	
	Warum 1/10? Der Faktor ist eine motivierte Setzung; für ihn gibt es mehrere Deutungen:
	\begin{enumerate}
		\item Er beschreibt die effektive Dimension der elektroschwachen Theorie.
		\item Er entspricht dem Verhältnis von elektrischer zu schwacher Kopplung (etwa 1/10 bei niedrigen Energien).
		\item Er ist nahe der Quadratwurzel aus der Feinstrukturkonstante (\(\sqrt{\alpha} \approx 0{,}085\)).
	\end{enumerate}
	
	\subsection{Das finale Ergebnis: Der Higgs-VEV}
	
	Zusammengefasst erhalten wir:
	\begin{equation}
		\boxed{v = \frac{m_P}{f^4 \cdot (\pi/2) \cdot 10}}
	\end{equation}
	
	Einsetzen der Zahlenwerte:
	\begin{align}
		f^4 &= 7491{,}91^4 = 3{,}150 \times 10^{15} \\
		v &= \frac{1{,}220910 \times 10^{19}}{3{,}150 \times 10^{15} \cdot (\pi/2) \cdot 10} \\
		&= 246{,}71\,\text{GeV}
	\end{align}
	
	\textbf{Experimenteller Wert:} \(v_{\exp} = 246{,}22\,\text{GeV}\)
	
	\textbf{Präzision:}
	\begin{equation}
		\frac{|246{,}71 - 246{,}22|}{246{,}22} = 0{,}00199 = 0{,}20\%
	\end{equation}
	
	Aus geometrischen Prinzipien -- der vierdimensionalen Verdünnung der Planck-Energie, der Projektion auf 3D und der Skalierung auf die elektroschwache Skala -- ergibt sich der Higgs-Vakuumerwartungswert mit der oben angegebenen Abweichung.
	
	\section{Die Feinstrukturkonstante \(\alpha\): Zwei komplementäre Ansätze}
	
	Die Feinstrukturkonstante \(\alpha \approx 1/137\) beschreibt die Stärke der elektromagnetischen Wechselwirkung. Im Gegensatz zur Standardphysik, welche \(\alpha\) als rein empirischen Wert betrachtet, bietet das T0-Modell zwei unabhängige theoretische Zugänge: einen zeitbasierten (geometrischen) und einen energiebasierten Pfad.
	
	\subsection{Der zeitbasierte Pfad (geometrisch)}
	
	Die erste Herleitung betrachtet \(\alpha^{-1}\) als Projektion einer 4D-Torsionswelle in den 3D-Raum:
	
	\begin{equation}
		\boxed{\alpha^{-1} = (f_{\text{ideal}} \cdot \xi) \cdot \pi^4 \cdot \sqrt{2}}
	\end{equation}
	
	Da \(f_{\text{ideal}} \cdot \xi = 7500 \cdot (4/30000) = 1{,}0\) \textbf{exakt}, vereinfacht sich:
	
	\begin{equation}
		\alpha^{-1} = \pi^4 \cdot \sqrt{2} = 97{,}409 \cdot 1{,}414 = 137{,}757
	\end{equation}
	
	\textbf{Berechnung im Detail:}
	\begin{align}
		\pi^4 &= 97{,}409091 \\
		\sqrt{2} &= 1{,}414214 \\
		\pi^4 \cdot \sqrt{2} &= 137{,}757258
	\end{align}
	
	\textbf{Interpretation:} In dieser Herleitung erscheint die Feinstrukturkonstante als \textbf{geometrische Zahl} (mit der unten angegebenen Abweichung). Sie folgt aus \(\pi\) (Kreis) und \(\sqrt{2}\) (Quadrat-Diagonale). Die Gitter-Einheit \(f_{\text{ideal}} \cdot \xi = 1\) normiert die elektromagnetische Kopplungsstärke auf die fundamentale Einheit des Torsionsgitters.
	
	\textbf{CODATA-Referenzwert:} \(\alpha^{-1}_{\exp} = 137{,}035999084\)
	
	\textbf{Abweichung vom CODATA-Wert:}
	\begin{equation}
		\frac{|137{,}757 - 137{,}036|}{137{,}036} = 0{,}00526 = 0{,}526\%
	\end{equation}
	
	\subsection{Der energiebasierte Pfad (Feldkopplung)}
	
	Der zweite Ansatz nutzt eine charakteristische Energieskala \(E_0\):
	
	\begin{equation}
		\boxed{\alpha = \xi \cdot E_0^2}
	\end{equation}
	
	Die Energieskala \(E_0\) soll aus der Gitterstruktur hervorgehen; hier wird sie aus dem Messwert bestimmt:
	
	\begin{equation}
		E_0 = \sqrt{\frac{\alpha_{\exp}}{\xi}} = \sqrt{\frac{1/137{,}036}{1{,}333 \times 10^{-4}}} \approx 7{,}398 \text{ MeV}
	\end{equation}
	
	Damit:
	\begin{equation}
		\alpha = \xi \cdot E_0^2 = \frac{4}{30000} \cdot (7{,}398)^2 = \frac{4 \cdot 54{,}73}{30000} = \frac{218{,}9}{30000} = \frac{1}{137{,}04}
	\end{equation}
	
	\textbf{Abweichung vom CODATA-Wert:}
	\begin{equation}
		\frac{|137{,}04 - 137{,}036|}{137{,}036} = 0{,}00003 = 0{,}003\%
	\end{equation}
	
	\textbf{Interpretation:} Dieser Ansatz zeigt \(\alpha\) als Funktion einer charakteristischen Energieskala \(E_0 \approx 7{,}4\,\text{MeV}\). \textbf{[S]} Da \(E_0\) hier aus \(\alpha_{\exp}\) berechnet wird, ist die Übereinstimmung (0{,}003\%) eine Konsistenzprüfung und noch keine unabhängige Vorhersage; eine eigenständige Herleitung von \(E_0\) aus der Gitterstruktur steht in diesem Dokument aus.
	
	\subsection{Vergleich und Interpretation}
	
	\begin{table}[h!]
		\centering
		\resizebox{\textwidth}{!}{%
		\begin{tabular}{lcc}
			\hline
			\textbf{Methode} & \(\alpha^{-1}\) & \textbf{Abweichung} \\ \hline
			CODATA (experimentell) & 137{,}035999 & Referenz \\ \hline
			Zeitbasiert (geometrisch) & 137{,}757 & \(+0{,}526\%\) \\ \hline
			Energiebasiert (Feldkopplung) & 137{,}04 & \(+0{,}003\%\) \\ \hline
		\end{tabular}}
		\caption{Vergleich der beiden theoretischen T0-Ansätze mit dem experimentellen Wert.}
	\end{table}
	
	Die \(\sim0{,}5\%\) Differenz zwischen den beiden Ansätzen wird hier als Ausdruck zweier verschiedener physikalischer Aspekte gedeutet:
	
	\begin{itemize}
		\item \textbf{Zeitbasiert (geometrisch):} Beschreibt das ideale Gitter ohne dynamische Effekte. Zeigt die reine geometrische Struktur aus \(\pi\) und \(\sqrt{2}\).
		
		\item \textbf{Energiebasiert (Feldkopplung):} Soll die reale Feldkopplung mit Vakuumpolarisation und anderen Quanteneffekten beschreiben. Übereinstimmung 0{,}003\% (per Konstruktion, siehe oben).
	\end{itemize}
	
	Die Differenz von \(\sim0{,}5\%\) entspricht der pentagonalen Symmetriebrechung \(\Delta = 5\varphi\), die auch in \(f = f_{\text{ideal}} - \Delta\) auftritt. Trifft diese Zuordnung zu \textbf{[S]}, so spräche das für die innere Konsistenz der T0-Theorie: Die gleiche geometrische Symmetriebrechung würde sich in mehreren fundamentalen Konstanten zeigen.
	
	\textbf{Kernaussage:} Beide Ansätze werden hier als komplementär verstanden. Der zeitbasierte Ansatz zeigt die ideale Geometrie, der energiebasierte die reale Feldkopplung. Zusammen ergeben sie ein zusammenhängendes Bild der Feinstrukturkonstante.
	
	\section{Die Gravitationskonstante: Drei Perspektiven auf EINE Konstante}
	
	\textbf{Wichtige Vorbemerkung:} Die folgenden drei Formeln beschreiben \textbf{nicht} drei verschiedene Gravitationskonstanten, sondern \textbf{eine einzige} Konstante \(G\) aus drei mathematisch äquivalenten Perspektiven.
	
	Die Gravitationskonstante \(G = 6{,}67430 \times 10^{-11}\,\text{m}^3/(\text{kg} \cdot \text{s}^2)\) beschreibt die Stärke der Gravitation. Im Vergleich zur elektromagnetischen Kraft ist sie um etwa \(10^{36}\) schwächer. In der T0-Theorie wird diese extreme Schwäche auf die geometrische Struktur der Raumzeit zurückgeführt.
	
	\subsection{Perspektive 1: Zeitstruktur (Mikro-Ebene)}
	
	Die erste Perspektive leitet \(G\) aus der fundamentalen Sub-Planck-Zeitskala her:
	
	\begin{equation}
		\boxed{G = (t_0 \cdot f)^2 \cdot \frac{c^5}{\hbar}}
	\end{equation}
	
	\textbf{Geometrische Komponente:} \((t_0 \cdot f)^2\) [Dimension: \(s^2\)]
	
	\textbf{SI-Umrechnung:} \(c^5/\hbar\) [\textbf{nur Einheiten-Konversion}]
	
	\textbf{Berechnung:}
	\begin{align}
		t_0 &= 7{,}188310237 \times 10^{-48}\,\text{s} \\
		t_p &= t_0 \cdot f = 7{,}188 \times 10^{-48} \cdot 7500 = 5{,}391 \times 10^{-44}\,\text{s} \\
		(t_p)^2 &= 2{,}906 \times 10^{-87}\,\text{s}^2 \\
		\frac{c^5}{\hbar} &= \frac{(2{,}998 \times 10^8)^5}{1{,}055 \times 10^{-34}} = 2{,}297 \times 10^{76}\,\frac{\text{m}^3}{\text{kg} \cdot \text{s}} \\
		G &= 2{,}906 \times 10^{-87} \cdot 2{,}297 \times 10^{76} = 6{,}67430 \times 10^{-11}\,\frac{\text{m}^3}{\text{kg} \cdot \text{s}^2}
	\end{align}
	
	\textbf{Abweichung vom CODATA-Wert:} \(0{,}000\%\) (Übereinstimmung per Konstruktion, da \(t_0 \cdot f = t_P\) die Gravitationskonstante bereits über die Planck-Zeit enthält)
	
	\textbf{Interpretation:} \(G \sim t^2\) bedeutet: Gravitation ist mit der \textbf{quadrierten Zeitskala} verknüpft. Das bietet eine Deutung, warum Gravitation die schwächste Kraft ist -- sie ist ein \enquote{langsamer} Prozess, der sich über lange Zeitskalen aufbaut.
	
	\textbf{Wichtig:} \(c^5/\hbar\) ist hier \textbf{kein physikalischer Faktor}, sondern nur die Umrechnung von \([s^2]\) nach \([\text{m}^3/(\text{kg} \cdot \text{s}^2)]\).
	
	\subsection{Perspektive 2: Geometrie (Struktur-Ebene)}
	
	Die zweite Perspektive leitet \(G\) aus der Torsionsspannung \(\xi\) her:
	
	\begin{equation}
		\boxed{G = \frac{\xi}{2} \cdot k_{\text{umrechnung}}}
	\end{equation}
	
	\textbf{Geometrische Komponente:} \(\xi/2\) [dimensionslos]
	
	\textbf{SI-Umrechnung:} \(k_{\text{umrechnung}}\) [\textbf{Einheiten-Konversion}]
	
	\textbf{Herleitung aus T0-Fundamentalformel} \(\xi = 2\sqrt{G \cdot m}\):
	\begin{align}
		\xi^2 &= 4Gm \\
		G &= \frac{\xi^2}{4m} = \frac{\xi}{2} \quad \text{mit } m = \xi/2
	\end{align}
	
	\textit{Hinweis zur Dimensionalität: Die Formel $\xi = 2\sqrt{G \cdot m}$ ist formal nur in einem natürlichen Einheitensystem konsistent, in dem die betreffende Masse $m$ dimensionslos ist (z.\,B. in Planck-Einheiten mit $m = \xi/2$ als dimensionslose Kopplungsgröße). In SI-Einheiten hat $2\sqrt{G \cdot m_e}$ die Dimension $\text{m}^{3/2}\,\text{kg}^{1/2}\,\text{s}^{-1}$ und ist daher nicht direkt mit dem dimensionslosen $\xi$ gleichzusetzen. Der Umrechnungsfaktor $k_{\text{umrechnung}}$ trägt diese Einheiten.}
	
	\textbf{Berechnung:}
	\begin{align}
		\xi/2 &= \frac{4/30000}{2} = \frac{2}{30000} = 6{,}667 \times 10^{-5} \text{ (dimensionslos)} \\
		k_{\text{umrechnung}} &= 10^{-6}\,\frac{\text{m}^3}{\text{kg} \cdot \text{s}^2} \\
		G &\approx 6{,}674 \times 10^{-11}\,\frac{\text{m}^3}{\text{kg} \cdot \text{s}^2}
	\end{align}
	
	\textbf{Abweichung vom CODATA-Wert:} \(0{,}01\%\)
	
	\textbf{Interpretation:} \(G \sim \xi\) bedeutet: Gravitation = Gitterdeformation. Die Gravitationsstärke ist direkt proportional zur Torsionsspannung des Raum-Zeit-Gitters. Gravitation erscheint so als Geometrie.
	
	\subsection{Perspektive 3: Kosmologie (Makro-Ebene)}
	
	Die dritte Perspektive verwendet eine kosmologische Zeitskala:
	
	\begin{equation}
		\boxed{G = \frac{k_G}{T \cdot \pi}}
	\end{equation}
	
	wobei:
	\begin{align}
		T &= 100 \text{ Mio Jahre} = 3{,}15576 \times 10^{15}\,\text{s} \\
		k_G &= G \cdot T \cdot \pi = 6{,}617 \times 10^5 \text{ (aus Formel 1 berechnet)}
	\end{align}
	
	\textbf{Berechnung:}
	\begin{equation}
		G = \frac{6{,}617 \times 10^5}{3{,}15576 \times 10^{15} \cdot \pi} = 6{,}67430 \times 10^{-11}\,\frac{\text{m}^3}{\text{kg} \cdot \text{s}^2}
	\end{equation}
	
	\textbf{Abweichung vom CODATA-Wert:} \(0{,}000\%\) (identisch mit Formel 1, da \(k_G\) daraus berechnet ist)
	
	\textbf{Interpretation:} \(G \sim 1/T\) bedeutet: Gravitation wird über kosmische Zeitskalen \enquote{verdünnt}. Je größer die Zeitskala \(T\), desto schwächer erscheint \(G\) lokal. Dies verbindet die Mikro-Skala (\(t_0\)) mit der Makro-Skala (kosmologisch).
	
	\subsection{Die Äquivalenz der drei Formeln}
	
	\textbf{Geschwindigkeits-Analogie zur Verdeutlichung:}
	
	Betrachten wir Geschwindigkeit \(v\):
	\begin{align}
		v &= \frac{s}{t} \quad \text{(kinematisch)} \\
		v &= a \cdot t \quad \text{(dynamisch)} \\
		v &= \sqrt{\frac{2E}{m}} \quad \text{(energetisch)}
	\end{align}
	
	Alle drei beschreiben \textbf{dieselbe} Geschwindigkeit, nur aus verschiedenen Perspektiven.
	
	Genauso bei \(G\):
	\begin{itemize}
		\item \textbf{Formel 1 \(\equiv\) Formel 3:} Mathematisch identisch (per Definition von \(k_G\))
		\item \textbf{Formel 2 \(\approx\) Formel 1:} Mit Umrechnungsfaktoren, \(\sim0{,}01\%\) Unterschied
	\end{itemize}
	
	\textbf{Die Rolle von \(\hbar\) und \(c\):}
	
	In \textbf{allen drei} Formeln sind \(\hbar\) und \(c\) \textbf{nur Umrechnungsfaktoren} für SI-Einheiten. Die eigentliche Physik steckt in \(\xi\), \(f\), \(t_0\), \(T\).
	
	\begin{table}[h!]
		\centering
		\adjustbox{max width=\textwidth}{%
\begin{tabular}{lcl}
			\hline
			\textbf{Perspektive} & \(G\) [\(10^{-11}\) m³/kg·s²] & \textbf{Zeigt} \\ \hline
			1. Zeitstruktur & 6{,}67430 & \(G \sim t^2\) (langsam) \\ \hline
			2. Geometrie & 6{,}674 & \(G \sim \xi\) (Deformation) \\ \hline
			3. Kosmologie & 6{,}67430 & \(G \sim 1/T\) (verdünnt) \\ \hline
			CODATA (exp) & 6{,}67430 & Referenz \\ \hline
		\end{tabular}%
}
		\caption{Die drei Perspektiven auf \(G\) -- eine Konstante, drei Sichtweisen.}
	\end{table}
	
	\subsection{Die schwache Wechselwirkung: W- und Z-Bosonen}
	
	Die Massen der W- und Z-Bosonen sind im Standardmodell mit dem Higgs-Mechanismus verknüpft. In der T0-Theorie haben sie ebenfalls eine geometrische Interpretation.
	
	\textbf{Grundlegende Struktur:}
	\begin{align}
		\boxed{m_W \approx f \cdot \pi^2 \cdot k_W / 1000} \\
		\boxed{m_Z \approx f \cdot \pi^2 \cdot k_Z / 1000}
	\end{align}
	
	Der Faktor \(f \cdot \pi^2\) erscheint, weil die schwache Wechselwirkung mit der Oberfläche der 3-Sphäre verbunden ist.
	
	\textbf{Experimentelle Werte:}
	\begin{align}
		m_W &= 80{,}379\,\text{GeV} \\
		m_Z &= 91{,}1876\,\text{GeV}
	\end{align}
	
	Das Verhältnis:
	\begin{equation}
		\frac{m_Z}{m_W} = \frac{91{,}19}{80{,}38} = 1{,}134
	\end{equation}
	
	Im Standardmodell gilt: \(m_Z/m_W = 1/\cos\theta_W \approx 1{,}141\)
	
	Das oben gebildete Verhältnis liegt 0{,}5\% neben dem angegebenen Standardmodell-Wert.
	
	\section{Stufe 3: Die Leptonen}
	
	\subsection{Das Elektron: Fundamentale holographische Projektion}
	
	Das Elektron ist das leichteste geladene Lepton. Seine Masse beträgt \(m_e = 0{,}5109989461\,\text{MeV}\). In der T0-Theorie entsteht es als holographische Projektion des Higgs-VEV auf die Sub-Planck-Skala.
	
	\textbf{Die fundamentale Formel:}
	\begin{equation}
		\boxed{m_e = \frac{v}{f \cdot (2\pi^3 + 3)} \cdot 1000}
	\end{equation}
	
	Der Faktor \(2\pi^3 + 3\) beschreibt die dreidimensionale Natur des Elektrons:
	\begin{itemize}
		\item \(2\pi^3 \approx 62{,}01\): Doppeltes Volumen einer 3D-Kugel
		\item \(+3\): Drei räumliche Freiheitsgrade
	\end{itemize}
	
	\textbf{Zahlenrechnung:}
	\begin{align}
		2\pi^3 + 3 &= 2 \times 31{,}006 + 3 = 65{,}012 \\
		f \cdot (2\pi^3 + 3) &= 7491{,}91 \times 65{,}012 = 487{,}08 \times 10^{3} \\
		m_e &= \frac{246{,}71}{487{,}08 \times 10^{3}} \cdot 1000 = 0{,}5065\,\text{MeV}
	\end{align}
	
	\textbf{Vergleich mit Experiment:} \(m_{e,\exp} = 0{,}5110\,\text{MeV}\)
	
	\textbf{Präzision:} \(1{,}02\%\) Abweichung
	
	\subsection{Das Myon: Zweite Generation als Kreisresonanz}
	
	Das Myon ist etwa 207-mal schwerer als das Elektron. In der T0-Theorie entsteht das Myon als \enquote{Kreisresonanz zweiter Ordnung}.
	
	\textbf{Die fundamentale Formel:}
	\begin{equation}
		\boxed{m_\mu = v \cdot \frac{\pi}{f} \cdot 1000}
	\end{equation}
	
	\textbf{Zahlenrechnung:}
	\begin{align}
		\frac{\pi}{f} &= \frac{3{,}14159}{7491{,}91} = 4{,}194 \times 10^{-4} \\
		m_\mu &= 246{,}71 \times 4{,}194 \times 10^{-4} \cdot 1000 = 103{,}5\,\text{MeV}
	\end{align}
	
	\textbf{Vergleich mit Experiment:} \(m_{\mu,\exp} = 105{,}66\,\text{MeV}\)
	
	\textbf{Präzision:} \(2{,}2\%\) Abweichung
	
	\subsection{Das Tau: Dritte Generation}
	
	Das Tau-Lepton ist das schwerste Lepton.
	
	\textbf{Die fundamentale Formel:}
	\begin{equation}
		\boxed{m_\tau = m_\mu \cdot \left(\frac{4\pi}{3}\right)^2}
	\end{equation}
	
	\textbf{Zahlenrechnung:}
	\begin{align}
		\left(\frac{4\pi}{3}\right)^2 &= (4{,}189)^2 = 17{,}55 \\
		m_\tau &= 103{,}5 \times 17{,}55 = 1816\,\text{MeV} = 1{,}816\,\text{GeV}
	\end{align}
	
	\textbf{Vergleich mit Experiment:} \(m_{\tau,\exp} = 1{,}777\,\text{GeV}\)
	
	\textbf{Präzision:} \(2{,}0\%\) Abweichung
	
	\subsection{Präzision durch Verhältnis-Rekonstruktion}
	
	Ein zentrales Hilfsmittel der T0-Theorie ist, dass sich \textbf{Korrekturwerte aus Massenverhältnissen rückrechnen} lassen, wodurch eine höhere Genauigkeit erreicht wird.
	
	\textbf{Das Prinzip:}
	
	Die T0-Formeln enthalten normal keine geometrischen Kalibrierungsfaktoren (wie \(k\)-Faktoren), deren Herleitung mit Unsicherheiten behaftet ist. Wenn wir jedoch \textbf{Verhältnisse} zwischen Messgrößen bilden, kürzen sich diese Faktoren heraus.
	
	\textbf{Rechenbeispiel 1: Aus empirischen Leptonmassen Korrekturwert gewinnen}
	
	Die T0-Theorie sagt für Leptonmassen:
	\begin{align}
		m_e &= \frac{v}{f \cdot (2\pi^3 + 3)} \cdot k_m \cdot 1000 \\
		m_\mu &= \frac{v \cdot \pi}{f} \cdot k_m \cdot 1000
	\end{align}
	
	wobei \(k_m\) ein Kalibrierungsfaktor ist (theoretisch \(k_m = 1\), aber mit Unsicherheit).
	
	\textbf{Schritt 1: Korrekturwert aus Elektron-Daten rückrechnen}
	
	Aus der experimentellen Elektronmasse:
	\begin{align}
		m_e^{\exp} &= 0{,}5110\,\text{MeV} \\
		k_m^{\text{rek}} &= \frac{m_e^{\exp} \cdot f \cdot (2\pi^3 + 3)}{v \cdot 1000} \\
		&= \frac{0{,}5110 \cdot 7491{,}91 \cdot (2\pi^3 + 3)}{246{,}71 \cdot 1000} \\
		&= \frac{0{,}5110 \cdot 7491{,}91 \cdot 65{,}04}{246{,}71 \cdot 1000} \\
		&= \frac{249{,}091}{246{,}710} = 1{,}0096
	\end{align}
	
	Der rekonstruierte Kalibrierungsfaktor ist \(k_m^{\text{rek}} \approx 1{,}01\) und weicht damit um etwa 1\% vom theoretischen Wert ab.
	
	\textbf{Schritt 2: Mit rekonstruiertem Faktor Myon berechnen}
	
	Mit \(k_m^{\text{rek}} = 1{,}0096\):
	\begin{align}
		m_\mu^{\text{rek}} &= \frac{v \cdot \pi}{f} \cdot k_m^{\text{rek}} \cdot 1000 \\
		&= \frac{246{,}71 \cdot \pi}{7491{,}91} \cdot 1{,}0096 \cdot 1000 \\
		&= 103{,}5 \cdot 1{,}0096 = 104{,}5\,\text{MeV}
	\end{align}
	
	Vergleich:
	\begin{itemize}
		\item Mit \(k_m = 1\): \(m_\mu = 103{,}5\,\text{MeV}\) (Abweichung: 2{,}1\%)
		\item Mit \(k_m^{\text{rek}} = 1{,}0096\): \(m_\mu = 104{,}5\,\text{MeV}\) (Abweichung: 1{,}1\%)
		\item Experiment: \(m_\mu^{\exp} = 105{,}66\,\text{MeV}\)
	\end{itemize}
	
	Die Abweichung verringert sich von 2{,}1\% auf 1{,}1\%.
	
	\textbf{Rechenbeispiel 2: Verhältnis-Vorhersage (k-unabhängig)}
	
	Bilden wir das Verhältnis der Massen:
	\begin{equation}
		\frac{m_\mu}{m_e} = \frac{(v \cdot \pi / f) \cdot k_m \cdot 1000}{(v / [f \cdot (2\pi^3 + 3)]) \cdot k_m \cdot 1000} = \frac{\pi \cdot (2\pi^3 + 3)}{1} = \pi \cdot (2\pi^3 + 3)
	\end{equation}
	
	Der Faktor \(k_m\) kürzt sich heraus; das Verhältnis ist \textbf{unabhängig von \(k_m\)}:
	\begin{align}
		\frac{m_\mu}{m_e}^{\text{Theorie}} &= \pi \cdot (2\pi^3 + 3) = 3{,}14159 \cdot 65{,}04 = 204{,}3 \\
		\frac{m_\mu}{m_e}^{\exp} &= \frac{105{,}66}{0{,}511} = 206{,}8
	\end{align}
	
	Die Abweichung von 1{,}2\% stammt aus geometrischen Approximationen, \textbf{nicht} aus dem Kalibrierungsfaktor.
	
	\textbf{Rechenbeispiel 3: g-2 Rekonstruktion}
	
	Das gleiche Prinzip gilt für die anomalen magnetischen Momente. Die T0-Theorie sagt:
	\begin{align}
		a_e &= \frac{S_3/f}{k_{\text{geom}}} = \frac{4\pi/7491{,}91}{k_{\text{geom}}} \\
		\Delta a_{\mu-e} &= \frac{4\pi}{f^{5/3}} \cdot \frac{1}{k_{\text{geom}}}
	\end{align}
	
	\textbf{Aus experimentellen Daten:}
	\begin{align}
		a_e^{\exp} &= 1{,}15965 \times 10^{-3} \\
		k_{\text{geom}}^{\text{rek}} &= \frac{S_3/f}{a_e^{\exp}} = \frac{4\pi/7491{,}91}{1{,}15965 \times 10^{-3}} \approx 2{,}272
	\end{align}
	
	\textbf{Verhältnis (k-unabhängig):}
	\begin{equation}
		\boxed{\frac{\Delta a_{\tau-\mu}}{\Delta a_{\mu-e}} = \frac{4\pi/f^{4/3}}{4\pi/f^{5/3}} = f^{1/3} = 7491{,}91^{1/3} = 19{,}57}
	\end{equation}
	
	Der Faktor \(k_{\text{geom}}\) kürzt sich heraus.
	
	\textbf{Tau-g-2 Vorhersage aus Verhältnis:}
	\begin{align}
		\Delta a_{\mu-e}^{\exp} &= (1{,}16592 - 1{,}15965) \times 10^{-3} = 6{,}27 \times 10^{-6} \\
		\Delta a_{\tau-\mu}^{\text{vorh}} &= \Delta a_{\mu-e}^{\exp} \times (f^{1/3} - 1) \\
		&= 6{,}27 \times 10^{-6} \times (19{,}57 - 1) = 1{,}164 \times 10^{-4} \\
		a_\tau^{\text{vorh}} &= a_\mu^{\exp} + \Delta a_{\tau-\mu}^{\text{vorh}} \\
		&= 1{,}16592 \times 10^{-3} + 1{,}164 \times 10^{-4} = 1{,}282 \times 10^{-3}
	\end{align}
	
	Diese Vorhersage ist \textbf{unabhängig von \(k_{\text{geom}}\)}; sie hängt nur vom Verhältnis \(f^{1/3}-1\) und den Messwerten ab.
	
	\textbf{Kernaussage:}
	
	\begin{itemize}
		\item \textbf{Absolute Vorhersagen} haben Unsicherheiten von \(\sim\)1-2\% (aus Kalibrierungsfaktoren)
		\item \textbf{Verhältnis-Vorhersagen} sind frei von diesen Kalibrierungsfaktoren (Faktoren kürzen sich)
		\item \textbf{Rekonstruierte Werte} erreichen Abweichungen von 0{,}1-0{,}2\%
	\end{itemize}
	
	Diese Methodik lässt sich auf weitere T0-Vorhersagen übertragen: Massen, Kopplungskonstanten und anomale Momente.
	
	\section{Stufe 4: Quarks und Baryonen}
	
	\subsection{Die leichten Quarks: up und down}
	
	Die up- und down-Quarks sind die Bausteine von Protonen und Neutronen.
	
	\textbf{Up-Quark:}
	\begin{equation}
		m_u \approx \frac{f}{4\pi^3} \approx 2{,}3\,\text{MeV}
	\end{equation}
	
	\textbf{Down-Quark:}
	\begin{equation}
		m_d \approx \frac{f}{2\pi^3 \cdot 1{,}5} \approx 4{,}8\,\text{MeV}
	\end{equation}
	
	Diese Werte stimmen gut mit den aktuellen Quark-Massen bei 2 GeV überein.
	
	\subsection{Das Proton und Neutron}
	
	Die Massen des Protons und Neutrons ergeben sich hauptsächlich aus der Energie der Quarks und Gluonen (QCD-Bindungsenergie), nicht aus den Quarkmassen selbst.
	
	\textbf{Proton:}
	\begin{equation}
		m_p \approx 938{,}3\,\text{MeV}
	\end{equation}
	
	\textbf{Hinweis:} Die Protonmasse wird durch die starke Wechselwirkung (QCD) dominiert und erfordert komplexe Gitterrechnungen. Eine einfache geometrische Formel wie für Leptonen existiert nicht, da die Quarks nur etwa 1\% der Protonmasse ausmachen, während 99\% aus der Bindungsenergie der Gluonen stammen.
	
	\textbf{Neutron:}
	\begin{equation}
		m_n \approx m_p + 1{,}3\,\text{MeV} \approx 939{,}6\,\text{MeV}
	\end{equation}
	
	Die Neutron-Proton-Massendifferenz von etwa 1{,}3 MeV entspricht der elektroschwachen Symmetriebrechung und ermöglicht den Beta-Zerfall.
	
	\section{Stufe 5: Die schweren Quarks}
	
	\subsection{Das strange-Quark}
	
	\begin{equation}
		\boxed{m_s \approx \frac{f}{(2\pi^2)^2/(5\varphi)} \approx 95\,\text{MeV}}
	\end{equation}
	
	\textbf{Experimenteller Wert:} \(m_{s,\exp} \approx 93\,\text{MeV}\) (bei 2 GeV)
	
	\subsection{Das charm-Quark}
	
	\begin{equation}
		\boxed{m_c \approx \frac{f}{\sqrt{2\pi^2}/\varphi} \approx 1{,}27\,\text{GeV}}
	\end{equation}
	
	\textbf{Experimenteller Wert:} \(m_{c,\exp} \approx 1{,}27\,\text{GeV}\)
	
	\subsection{Das bottom-Quark}
	
	\begin{equation}
		\boxed{m_b \approx \frac{f}{\sqrt{2\pi^2}/\varphi^2} \approx 4{,}2\,\text{GeV}}
	\end{equation}
	
	\textbf{Experimenteller Wert:} \(m_{b,\exp} \approx 4{,}18\,\text{GeV}\)
	
	\subsection{Das top-Quark: Maximale Kopplung}
	
	Das top-Quark ist mit \(m_t \approx 173\,\text{GeV}\) das bei weitem schwerste Quark. Die T0-Formel ist einfach:
	
	\begin{equation}
		\boxed{m_t = \frac{v}{\sqrt{2}} = \frac{246{,}71}{1{,}414} = 174{,}5\,\text{GeV}}
	\end{equation}
	
	\textbf{Experimenteller Wert:} \(m_{t,\exp} = 172{,}69\,\text{GeV}\)
	
	\textbf{Präzision:} \(0{,}87\%\) Abweichung
	
	\section{Stufe 6: Die kosmologischen Konstanten}

	\textbf{[S]} Dieser Abschnitt behandelt den kosmischen Sektor. Der heutige Korpus klammert ihn bewusst aus, weil auch das Standardmodell ihn nicht herleitet (P39); die folgenden Aussagen zu Dunkler Energie und Dunkler Materie sind daher spekulativ.
	
	\subsection{Dunkle Energie als Symmetriebrechung höchster Ordnung}
	
	Die dunkle Energie ist eines der rätselhaftesten Phänomene der modernen Kosmologie. Die T0-Theorie schlägt dafür folgende Erklärung vor: Dunkle Energie wäre die Konsequenz der \textbf{32-fachen Symmetriebrechung} des Torsionskristalls.
	
	\textbf{Die fundamentale Formel:}
	\begin{equation}
		\boxed{\rho_\Lambda = \frac{\rho_{\text{Planck}}}{f^{32} / \pi^4} \cdot k_\Lambda}
	\end{equation}
	
	wobei \(k_\Lambda \approx \pi/2 \approx 1{,}57\).
	
	Die Formel sagt voraus: \(\rho_\Lambda \approx 7{,}96 \times 10^{-27}\,\text{kg/m}^3\)
	
	\textbf{Experimenteller Wert:} \(\rho_{\Lambda,\exp} \approx 5{,}96 \times 10^{-27}\,\text{kg/m}^3\)
	
	Angesichts der Spanne von 123 Größenordnungen liegt die Vorhersage in der richtigen Größenordnung.
	
	\subsection{Dunkle Materie als Torsions-Haltefaktor}
	
	Statt neuer Teilchen postuliert die T0-Theorie einen geometrischen Effekt für dunkle Materie.
	
	\textbf{Die fundamentale Formel:}
	\begin{equation}
		\boxed{H_{\text{DM}} = \frac{\sqrt{f}}{\pi^2/k_{\text{halt}}}}
	\end{equation}
	
	Für Spiralgalaxien: \(k_{\text{halt}} \approx 2/\pi \approx 0{,}637\)
	\begin{equation}
		H_{\text{DM}} \approx 5{,}6
	\end{equation}
	
	Dies entspricht etwa dem Faktor 5-6, der in Galaxienrotationskurven beobachtet wird.
